\documentclass{article}
\usepackage{amsmath, amssymb, amsthm}
\usepackage{hyperref}
\usepackage{geometry}
\geometry{margin=1in}

\title{Erd\H{o}s Problem \#954}
\date{Last updated: 2025-08-31}
\author{Source: erdosproblems.com}

\begin{document}
\maketitle

\noindent\textbf{Status:} OPEN \quad \textbf{No prize}

\noindent\textbf{Tags:} number theory

\medskip
\noindent\textbf{Problem Statement:}

Let $0=a_0<a_1<a_2<\cdots$ be the sequence of integers defined by $a_0=0$ and $a_1=1$, and $a_{k+1}$ is the smallest integer $n$ for which the number of solutions to $a_i+a_j \leq n$ (with $0\leq i\leq j\leq k$ and $j\geq 1$) is $<n$.

Is the number of solutions to $a_i+a_j \leq x$ equal to $x+O(x^{1/4+o(1)})$?

\bigskip
\noindent\textbf{Remarks:}

This sequence was constructed by Rosen. Note that the number of solutions to $a_i+a_j\leq x$ is always at least $x$ by construction, but Erd\H{o}s and Rosen could not even prove whether it is at most $(1+o(1))x$.

The sequence begins\[0,1,3,5,9,13,17,24,31,38,45,\ldots.\]

\bigskip
\noindent\small{Source: \url{https://www.erdosproblems.com/954}}
\end{document}
